\section{Introduction}

Finite-wavenumber zero-group-velocity (ZGV) Lamb states are central to
long-lived localization in plates.  They retain a nonzero in-plane wavelength
while a stationary guided-wave dispersion suppresses the
leading transport of energy away from the source, producing local vibration
and a temporal decay law controlled by the geometry of the stationary set.
Laser-ultrasonic excitation and local ZGV responses are established
phenomena \citep{prada2005laser,prada2008local}.  In an isotropic plate,
rotational symmetry turns a radial finite-wavenumber stationary point into a
continuous ZGV critical ring, and radial stationary phase gives the established
\(t^{-1/2}\) decay for a nonnodal local response
\citep{prada2008power,laurent2014temporal}.  Neither the ZGV critical ring nor
this exponent is at issue here as a new effect; they provide the known
isotropic endpoint whose loss of directional degeneracy we examine.

Elastic anisotropy provides the other established endpoint.  Anisotropic-plate
studies established directional ZGV shifts \citep{prada2009anisotropy}.
Calculations also found multiple stationary points
\citep{hussain2012multiple,karous2019multiple}, and isolated anisotropic ZGV
points were previously reported.  More specifically, four minima and four saddles, together with
their interference, were previously reported for a cubic silicon plate
\citep{kiefer2023beating}; subsequent work described their directional
stationary-phase organization and unfolding \citep{kiefer2025skewing}.
Algorithms for computing isolated ZGV points in the full guided-wave spectrum
are likewise established \citep{kiefer2023computing,plestenjak2025sylvester}.
The present calculation therefore does not claim the isolated points, their
alternating four-plus-four organization, or the associated interference as
new observations.  Its question is how these known anisotropic objects emerge
in a controlled limit from the known isotropic ZGV critical ring.

A plate-specific question remains open in the literature examined here: how to
construct a coefficient-resolved connection between that geometry and the
temporal asymptotics of the full Lamb eigenproblem.
Morse--Bott critical manifolds and their perturbations are established
mathematics \citep{bott1954critical,banyaga2013cascades}, while
two-dimensional stationary phase has long been applied to isolated guided-wave
points \citep{velichko2007excitation,chapuis2010focusing,
karmazin2013caustics}.  Broken-symmetry uniform approximations in other wave
settings also connect continuous families to isolated objects through a
Bessel modulation \citep{creagh1996broken,brack1999uniform,brack2009closed}.
The required connection must retain full-elastic angular and radial
coefficients, a derivative-level remainder, source normalization, and a
phase-controlled temporal comparison.  This missing link, rather than a new
general Morse or Bessel theorem, is the focus of the analysis.

Here we connect the known ZGV critical ring to the anisotropic endpoint.  The
reduction treats weak cubic anisotropy on a simple tracked Lamb branch.
Full-elastic generalized-eigenvalue sensitivities determine the angular
coefficient \(V_4\) and the radial coefficient \(B\) without fitting.
Critical-point geometry, Cartesian Hessian inertia, and annular index
checks then test the alternating local Morse set and its perturbative
scaling inside a declared local annulus.  Uniform Bessel asymptotics connect
the ring response to the isolated-point regime, while the fixed-anisotropy
endpoint retains exact critical frequencies and Hessians.  Finally, validated
spectral computations and an accumulated inter-resolution phase-discrepancy
estimator test the entire chain without fitted amplitude, phase, frequency, or
time shifts.  The study is
computation-only: it concerns one selected simple branch and a controlled
weak-anisotropy family; the finite-anisotropy silicon stress test is separate
from the asymptotic evidence.
